\section{Introduction}
Generative search engines answer a query with a few sentences and attach URLs to them. The links make an answer look checkable, and answers are kept and reused: they are shared as links, pasted into reports and stored in chat histories. Whether the cited page supports a statement, however, is a property of the page \emph{at a given time}. The Web changes underneath citations: in scholarly articles, ``one in five articles suffers from reference rot''~\cite{Klein_2014} and ``three out of four URI references lead to changed content''~\cite{Jones_2016}. For AI search answers, verifiability has been measured once, at generation time, by human annotators~\cite{Liu_2023} or by automatic auditors~\cite{Wu_2025,venkit2025deeptraceauditingdeepresearch}.

Existing tools cover the two halves of the problem separately. Link checkers such as urlhealth~\cite{rao2026detectingcorrectingreferencehallucinations} decide whether a cited URL resolves, but a page that still loads can have stopped saying what the answer claims. Claim verifiers and attribution evaluators~\cite{Tang_2024,Gao_2023,Yue_2023} judge a statement against the document they are given; applied to a live citation, that is the page as it is today, so they cannot tell a citation that never supported a statement from one that did and then changed.

CiteDrift joins the two. For each statement--citation pair it judges support on the page as cited (an archived capture, or the page text saved with a dataset) and on the page today, and assigns one of six outcomes: \emph{still supported}, \emph{content changed}, \emph{uncertain}, \emph{link lost}, \emph{could not check}, or \emph{not supported when cited}. The contributions are:
(1) a two-time, claim-level citation checker with a calibrated link-state step that separates bot walls and network filters from real loss and judges redirects by content;
(2) a fully local implementation (no commercial API; one 12\,GB GPU) with a web explorer and a check mode for any answer; and
(3) a measurement: re-checking the human-annotated development and test citations of Liu et al.~\cite{Liu_2023} after about three and a half years, \linkLostRate of the checkable citations judged fully supporting in 2023 lost their link (a link-only checker reports \uhLiveLinkLostPct of these as live), and CiteDrift flags another \changedRate as content changes.
The demo relates to the Responsible Web track (information integrity, auditing), Search and Retrieval-Augmented AI (evaluation of AI search), and Web Mining (Web evolution).
